\begin{abstract}
In interactive imitation learning, a learner practises a task and an expert takes
over to demonstrate the right actions. Every takeover costs expert time, which
compute cannot reduce. Under a fixed budget, what each demonstration contains
decides the result. Existing methods decide only when to take over. They call the
expert at the first step of an episode where the policy becomes uncertain. An
episode that runs from start to end without achieving the task is a failed
episode. Two decisions are then left to chance: which failed episode to correct,
and where the demonstration begins.
We present \textbf{DISEIL} (\textbf{D}emonstration d\textbf{I}stillation
for \textbf{S}ample-\textbf{E}fficient \textbf{I}mitation \textbf{L}earning), which
decides both. After each round of practice, DISEIL groups the failed episodes into
modes of failure. For each mode it prescribes where a new demonstration should
begin, and checks the start is reachable before any expert time is spent. We
evaluate on five tasks, with image and state observations, ten settings in all,
at a budget of 20 demonstrations, against the best baselines. DISEIL reaches the
highest mean success rate in all ten, on average about 3 percentage points above
the strongest baseline.
\end{abstract}

\section{Introduction}

Behaviour cloning fits a policy to an expert's actions on the states the expert
visited \cite{pomerleau1988alvinn,bain1995cloning}. The learner's own small errors
then carry it into states it never saw, where errors compound. That drift is
called covariate shift
\cite{shimodaira2000covariate,ross2011dagger}. Interactive imitation learning
stops the drift: the learner acts, and the expert corrects it on the states the
learner actually reaches \cite{ross2011dagger,celemin2022iil}. The price is expert
effort, the cost that binds \cite{mandlekar2018roboturk,khazatsky2024droid}.
Policy performance scales with how much of the task the demonstrations cover, and
extra demonstrations inside a covered region add little
\cite{lin2024datascaling}. So once an allowance of $B$ demonstrations is fixed,
what each one contains is all that is left to choose (Figure~\ref{fig:teaser}).

\begin{figure}[t]
\centering
\includegraphics[width=0.62\columnwidth]{figures/teaser.png}
\caption{A round of interactive imitation learning ends with a set of failures
and one demonstration to spend on it. Existing methods decide only when to ask.
DISEIL also decides which failure the call is spent on, and which configuration
the expert starts from.}
\label{fig:teaser}
\end{figure}

The DAgger family calls the expert only when a test on the current state fires.
Methods differ only in what that test measures. Some measure the gap between the
learner's action and the expert's \cite{zhang2017safedagger}. Others measure the
spread of predictions under dropout \cite{menda2017dropoutdagger} or across an
ensemble \cite{menda2019ensembledagger}. Others use novelty and risk
\cite{hoque2021thriftydagger}, or a diffusion policy's own per-step denoising loss
\cite{lee2025diffdagger,chi2023diffusionpolicy}. Each reduces the state to one
number and uses it the same way. Asking for a demonstration takes three decisions.
When to call the expert. Which of the round's failed attempts to spend the call
on. Where the corrective demonstration begins. Every method above decides when. Which
attempt and where to begin are then settled by default. The attempt corrected is
whichever one fired the test first. The demonstration begins in the state that
attempt was already in. Under an unbounded budget those two defaults cost nothing,
because every mode is corrected eventually. Under a fixed budget they are the
whole problem.

One number attached to one state cannot say whether two failures are the same
mistake or two different ones. So a mode that recurs often is corrected often,
while a rare mode waits. The demonstration also starts at the state that fired
the test, frequently one the policy has already ruined. The expert then shows how
to recover from the damage, not how to avoid it. The failures within a round form
a small set with structure, and geometry recovers it. We describe the
configuration of the robot and the objects at the step where the policy first
becomes uncertain. Failures grouped by that description separate into
geometrically distinct modes. Whether those modes are worth allocating a budget
over is an empirical question, and our ablations answer it.

Demonstration distillation for sample-efficient imitation learning (DISEIL) keeps
the existing test for when to call the expert. It then decides which failure to
correct and where the demonstration begins. Each round, DISEIL anchors every
failed attempt at the step where its loss first crosses the threshold, not at the
step of highest loss. It groups the failures by the configuration at that anchor
step. Among the modes close to the largest in size, it targets the one with the
highest mean peak loss. For that mode it prescribes a starting configuration not
yet demonstrated. The prescribed start is checked against a constraint graph of
what the task permits, before any expert time is spent. A round of failures is
distilled into one demonstration that stands for a mode, not for a single episode.

We contribute three things. First, we separate the rule for asking into three
decisions: when, which and where. The DAgger family answers only when. Second, we
present DISEIL. It reaches the highest mean success rate on unseen test episodes
in all ten settings, on average 2.80 percentage points above the strongest
baseline in each.
Third, we locate the advantage in how the budget is allocated, not in how
informative each demonstration is on its own. Removing the mode structure costs a
mean 4.37 points of success while information gain does not fall, the signature of
a redundantly spread budget.
